%=======================================================================
%  受迫耦合振子模型 —— 第 2 章 + 附录 A
%  编译： xelatex main  →  bibtex main  →  xelatex main  →  xelatex main
%  或：   latexmk -xelatex main.tex
%
%  说明：本文件为驱动文件。正文内容在
%          chapter2_body.tex   （第 2 章）
%          appendixA_body.tex  （附录 A）
%  可直接把这两个文件 \input 进你自己的手稿，无需本驱动文件。
%=======================================================================
\documentclass[11pt,a4paper]{report}

%-----------------------------------------------------------------------
% 中文支持：双路径。有 ctex 用 ctex（推荐，断行/标点更规范）；
% 否则回落到 fontspec + Noto CJK，保证最小环境也能编译。
%-----------------------------------------------------------------------
\newif\ifhasctex
\IfFileExists{ctex.sty}{\hasctextrue}{\hasctexfalse}

\ifhasctex
  \usepackage[UTF8]{ctex}
\else
  \usepackage{fontspec}
  \IfFontExistsTF{Noto Serif CJK SC}
    {\setmainfont{Noto Serif CJK SC}}
    {\IfFontExistsTF{WenQuanYi Zen Hei}
      {\setmainfont{WenQuanYi Zen Hei}}
      {}}
  % 最小环境下手工设置 CJK 断行惩罚，避免中文不换行溢出版面
  \XeTeXlinebreaklocale "zh"
  \XeTeXlinebreakskip = 0pt plus 1pt minus 0.1pt
\fi

%-----------------------------------------------------------------------
% 数学与排版
%-----------------------------------------------------------------------
\usepackage{amsmath,amssymb,amsthm,mathtools}
\usepackage{bm}
\usepackage{geometry}
\geometry{margin=2.6cm}
\usepackage{booktabs,tabularx,array,multirow}
\usepackage{enumitem}
\usepackage{xcolor}
\usepackage{caption}
\captionsetup{font=small,labelsep=quad}
\usepackage{microtype}
\usepackage{natbib}
\bibliographystyle{plainnat}
\usepackage[colorlinks=true,linkcolor=blue!55!black,citecolor=teal!60!black,
            urlcolor=magenta!60!black,breaklinks=true]{hyperref}
\usepackage[nameinlink]{cleveref}

%-----------------------------------------------------------------------
% 公式按章编号 → (2.1), (2.2), ... ; 附录 → (A.1), ...
%-----------------------------------------------------------------------
\numberwithin{equation}{chapter}
\allowdisplaybreaks[2]

%-----------------------------------------------------------------------
% cleveref 中文名称
%-----------------------------------------------------------------------
\crefname{equation}{式}{式}
\Crefname{equation}{式}{式}
\crefname{chapter}{第}{第}
\crefname{section}{}{}
\crefname{table}{表}{表}
\crefname{appendix}{附录}{附录}
\crefname{remark}{说明}{说明}
\Crefname{remark}{说明}{说明}
\crefname{hypo}{命题}{命题}
\Crefname{hypo}{命题}{命题}

%-----------------------------------------------------------------------
% 记号宏（集中定义，便于全局替换）
%-----------------------------------------------------------------------
\newcommand{\wbar}{\bar{\omega}}          % 个体 mu 峰频
\newcommand{\geff}{\gamma_{\mathrm{eff}}} % 有效展宽
\newcommand{\Feff}{F_{\mathrm{eff}}}
\newcommand{\Om}{\Omega}
\newcommand{\Ostar}{\Om^{*}}
\newcommand{\dd}{\mathrm{d}}
\newcommand{\ii}{\mathrm{i}}
\newcommand{\Var}{\operatorname{Var}}
\newcommand{\Rea}{\operatorname{Re}}
\newcommand{\Ima}{\operatorname{Im}}
\newcommand{\Kc}{K_{\mathrm{c}}}
\newcommand{\Thf}{\Theta_{\mathrm{eff}}}
% 以下两个宏在正文与数学模式中都会用到，故用 \ensuremath 包裹
\newcommand{\TPM}{\ensuremath{\mathrm{TPM}}}
\newcommand{\YGTSS}{\ensuremath{\mathrm{YGTSS}}}

% 强调框（用于核心公式）
\newcommand{\keyeq}[1]{\boxed{\,#1\,}}

% 命题环境
\newtheorem{hypo}{命题}
\newtheorem{remark}{说明}[chapter]
\theoremstyle{definition}

%-----------------------------------------------------------------------
\title{\bfseries 固定频率正中神经刺激的受迫耦合振子理论\\[4pt]
       \large ——从节律带动假设到个体化最优频率}
\author{}
\date{}

\begin{document}
\maketitle
\tableofcontents

% 使首个 \chapter 编号为 2
\setcounter{chapter}{1}

%=======================================================================
%  第 2 章  受迫耦合振子模型
%  可独立 \input 进任意手稿。依赖：amsmath, booktabs, tabularx,
%  cleveref, natbib，以及 main.tex 中定义的记号宏。
%=======================================================================
\chapter{受迫耦合振子模型：从带动假设到个体化最优频率}
\label{chap:model}

%-----------------------------------------------------------------------
\section{本章目标与逻辑路线}
\label{sec:roadmap}

第~1 章确立了经验事实：固定 $10\,$Hz 的节律性正中神经刺激（rhythmic median
nerve stimulation, rMNS）在双盲假刺激对照试验中显著降低抽动严重度——活跃组
Yale 全球抽动严重度量表总分（\YGTSS-TTSS）平均下降 $7.13$ 分，假刺激组
$2.13$ 分，等待组 $2.26$ 分，效应量 $0.5$，$p=0.02$；应答率 $59\%$
对 $33\%$~\citep{Maiquez2023}。但\emph{机制未定}：独立复现研究显示节律性与
非节律性刺激的疗效无显著差异，其作者明确结论为脑电带动假设无法解释该
疗效~\citep{Iverson2023a}。

本章不预设带动假设为真。本章的任务是：\textbf{把带动假设写成一个封闭的动力学
模型，推导它的必然推论，使其成为可判伪的对象}。逻辑路线为
%
\begin{equation*}
\text{生理约束}\;\longrightarrow\;\text{相位约化}\;\longrightarrow\;
\text{受迫 Kuramoto}\;\longrightarrow\;\text{OA 精确降维}\;\longrightarrow\;
\rho^{*}(\Om)\;\longrightarrow\;E(\Om)\;\longrightarrow\;\text{可判伪预测}.
\end{equation*}
%
每一步的近似条件与适用域都在该步给出；模型解释不了的现象集中列于
\cref{sec:boundary}。完整代数推导见\cref{app:derivation}。

%-----------------------------------------------------------------------
\section{三个可分离的命题}
\label{sec:hypotheses}

带动假设的口语表述（“$10\,$Hz 刺激拉动 mu 节律，增强皮层抑制，从而减少
抽动”）包含三个逻辑上独立、可分别检验的命题。

\begin{hypo}[锁定命题]\label{hyp:H1}
感觉运动 mu 节律可近似为一群弱耦合极限环振子的集体行为；rMNS 构成外部周期
驱动。集体锁定程度由驱动强度、群体异质性、内部耦合与\emph{去谐量}
$\Delta=\wbar-\Om$ 共同决定。
\end{hypo}

\begin{hypo}[门控命题]\label{hyp:H2}
抽动的瞬时发生率是 mu 相干程度的单调递减函数。
\end{hypo}

\begin{hypo}[个体化命题]\label{hyp:H3}
存在个体依赖的最优刺激频率 $\Ostar_i$；固定频率方案对
$|\wbar_i-\Om_{\mathrm{fix}}|$ 较大的个体系统性失效。
\end{hypo}

\noindent\textbf{关键结构.}
\cref{hyp:H3} 不是附加假设，而是 \cref{hyp:H1} 与 \cref{hyp:H2} 的\emph{演绎
推论}。这一点决定了本章的全部价值：
%
\begin{quote}
若带动假设成立，频率个体化不是可选的工程优化，而是逻辑必然。
反之，若实测疗效对刺激频率不敏感，则 \cref{hyp:H1} 与 \cref{hyp:H2}
中至少一者为假。
\end{quote}
%
因此对 \cref{hyp:H3} 的检验同时是对带动假设本身的检验。这是本模型区别于
“给既有疗法套一层数学包装”的根本之处。

%-----------------------------------------------------------------------
\section{相位约化的合法性}
\label{sec:reduction}

在写下方程之前必须交代三个约化步骤，否则模型的适用范围不清。

\paragraph{步骤 1：从神经元到振子群.}
感觉运动皮层 mu 节律（$8$–$13\,$Hz）是同步化突触后电位求和的宏观表现。将
参与该节律的局部神经元集团粗粒化为 $N$ 个“振子单元”，每单元有自身固有频率
$\omega_i$。代价是丢弃单神经元放电细节，保留集体相位信息——而
EEG/MEG 测量的恰恰是集体量。

\paragraph{步骤 2：从极限环到相位.}
若各单元处于稳定极限环上且驱动足够弱（不足以显著改变振幅），则可作相位约化，
状态由单一相位变量 $\theta_i$ 描述。\textbf{这是本模型最强的一个近似}，其代价
见 \cref{sec:boundary} 第 (2) 条（真实 mu 是爆发式而非持续振荡）。

\paragraph{步骤 3：耦合函数取一阶.}
相位耦合函数 $\Gamma(\theta_j-\theta_i)$ 取其傅里叶展开的一阶项
$\sin(\theta_j-\theta_i)$。这保证解析可解性，代价是\emph{只保留 $1{:}1$ 锁定}，
高阶 $p{:}q$ 锁定被排除（\cref{sec:pulse} 单独处理）。

三步都成立时，得到受迫 Kuramoto 模型；三步之外的现象不在本章解释范围内。

%-----------------------------------------------------------------------
\section{M0：单振子基线}
\label{sec:M0}

先解最小模型，目的有二：建立 Arnold tongue 的物理直觉，以及为数值代码提供
解析金标准。

\subsection{无噪声：Adler 方程}

单个固有频率 $\omega_0$ 的振子受频率 $\Om$、强度 $F$ 的正弦驱动：
$\dot\theta=\omega_0+F\sin(\Om t-\theta)$。引入相对相位
$\varphi=\theta-\Om t$ 与去谐 $\Delta=\omega_0-\Om$：
%
\begin{equation}\label{eq:adler}
\keyeq{\dot\varphi=\Delta-F\sin\varphi}
\end{equation}
%
不动点条件 $\sin\varphi^{*}=\Delta/F$ 有解的充要条件是
%
\begin{equation}\label{eq:lock}
|\Delta|\le F
\qquad\Longleftrightarrow\qquad
|\Om-\omega_0|\le F .
\end{equation}
%
稳定性要求 $\partial_\varphi\dot\varphi=-F\cos\varphi^{*}<0$，故稳定分支为
$\varphi^{*}=\arcsin(\Delta/F)\in(-\pi/2,\pi/2)$。

\cref{eq:lock} 即 $1{:}1$ \textbf{Arnold tongue}：在 $(\Om,F)$ 平面上是顶点位于
$(\omega_0,0)$、边界斜率 $\pm1$ 的楔形区域。物理含义直接：\emph{驱动越强，容许
的频率失配越大}。这正是经颅交流电刺激（tACS）文献中倒三角示意图的严格
来源~\citep{Ali2013,Notbohm2016}。

\subsection{加噪声：倾斜洗衣板}

\begin{equation}\label{eq:sde}
\dd\varphi=(\Delta-F\sin\varphi)\,\dd t+\sqrt{2D}\,\dd W ,
\end{equation}
%
对应势函数 $U(\varphi)=-\Delta\varphi-F\cos\varphi$（验证：
$-\partial_\varphi U=\Delta-F\sin\varphi$）。零去谐时有\emph{精确闭式解}：
$\Delta=0$ 时 $P(\varphi)\propto e^{(F/D)\cos\varphi}$，即浓度参数
$\kappa_{\mathrm{vM}}=F/D$ 的 von Mises 分布，锁相值
%
\begin{equation}\label{eq:vonmises}
\rho\;\equiv\;\bigl|\langle e^{\ii\varphi}\rangle\bigr|
=\frac{I_1(F/D)}{I_0(F/D)}
\;\xrightarrow[\;F\ll D\;]{}\;\frac{F}{2D}.
\end{equation}

\begin{remark}\label{rem:Dposition}
注意 \cref{eq:vonmises} 的结构：噪声 $D$ 在线性响应阶出现在分母，系数为
$1/2$。这个位置与 \cref{sec:regimes} 中异质性宽度 $\gamma$ 的位置\emph{完全
相同}——这为 \cref{sec:noise} 合并二者为 $\geff=\gamma+D$ 提供了第一条独立
依据。
\end{remark}

M0 在工程上的用途：\cref{eq:lock} 的边界与 \cref{eq:vonmises} 的
$I_1/I_0$ 曲线是代码单元测试的硬断言。

%-----------------------------------------------------------------------
\section{M1：受迫 Kuramoto 模型}
\label{sec:M1}

\subsection{微观方程}

\begin{equation}\label{eq:micro}
\dot\theta_i=\omega_i
+\frac{K}{N}\sum_{j=1}^{N}\sin(\theta_j-\theta_i)
+F\sin(\Om t-\theta_i)
+\sqrt{2D}\,\xi_i(t)
\end{equation}
%
各量的生理对应（这一步的明确性决定模型能否与数据对接）见 \cref{tab:params}。

\begin{table}[htbp]
\centering
\caption{模型参数与生理/实验量的对应}
\label{tab:params}
\small
\begin{tabularx}{\textwidth}{@{}l l X l@{}}
\toprule
符号 & 数学含义 & 生理／实验对应 & 经验来源 \\
\midrule
$\wbar$ & $g(\omega)$ 中心 & \textbf{个体 mu 峰频} & 静息 EEG (C3/CP3) + specparam \\
$\gamma$ & $g(\omega)$ 半宽 & mu 谱峰宽度 & 同上 \\
$K$ & 内部耦合 & 内源相干性／谱峰高度 & 同上 \\
$F$ & 外部驱动 & 刺激强度 (mA) 经募集曲线映射 & ITC–强度剂量曲线 \\
$D$ & 噪声强度 & 感觉运动噪声，\textbf{TS 中升高} & \citet{Gialopsou2025} \\
$\Om$ & 驱动频率 & \textbf{设备设定，当前固定 $10\,$Hz} & --- \\
\bottomrule
\end{tabularx}
\end{table}

固有频率取洛伦兹分布
%
\begin{equation}\label{eq:lorentz}
g(\omega)=\frac{\gamma}{\pi\bigl[(\omega-\wbar)^2+\gamma^2\bigr]} .
\end{equation}
%
选洛伦兹\emph{不是}因为它比高斯更真实，而是因为它使下一步降维\emph{精确}。
这一选择的代价与检验方式见 \cref{sec:boundary} 第 (3) 条。

\subsection{两个驱动可合并为一个：模型的核心结构}

定义序参量
%
\begin{equation}\label{eq:orderparam}
z=r e^{\ii\psi}=\frac{1}{N}\sum_{j=1}^{N}e^{\ii\theta_j} ,
\end{equation}
%
则 \cref{eq:micro} 中两个驱动项可合并：
%
\begin{equation*}
Kr\sin(\psi-\theta_i)+F\sin(\Om t-\theta_i)
=\Ima\!\left[\bigl(Kz+Fe^{\ii\Om t}\bigr)e^{-\ii\theta_i}\right],
\end{equation*}
%
即每个振子只感受到\emph{单一等效复驱动场}
%
\begin{equation}\label{eq:Hfield}
\keyeq{H=Kz+Fe^{\ii\Om t}}
\end{equation}

\noindent\textbf{这是整个模型的物理核心。} 它说明“互同步”与“受迫同步”不是两个
独立机制，而是同一个驱动场的两个分量，二者在竞争同一个相位。当 $K$ 大时群体
倾向于锁到自己的集体相位 $\psi$；当 $F$ 大时倾向于锁到外部相位 $\Om t$。
临床上这意味着：\emph{内源节律强的个体更“顽固”，更难被外部驱动改写}——这条
定性推论将在 \cref{sec:regimeI} 得到定量形式。

%-----------------------------------------------------------------------
\section{Ott--Antonsen 精确降维}
\label{sec:OA}

对形如 $\dot\theta_i=\omega_i+\Ima(He^{-\ii\theta_i})$ 的系统，在洛伦兹分布
\cref{eq:lorentz} 下、$N\to\infty$、$D=0$ 时，Ott--Antonsen (OA)
ansatz 精确成立~\citep{Ott2008}，序参量满足
%
\begin{equation}\label{eq:oa}
\dot z=(\ii\wbar-\gamma)z+\tfrac12\bigl(H-H^{*}z^{2}\bigr).
\end{equation}

\paragraph{自检.}
取 $H=Kz$（无外驱动），得
$\dot z=(\ii\wbar-\gamma)z+\frac{K}{2}z(1-|z|^2)$。在 $\wbar=0$、$z=r$ 实数时
$\dot r=-\gamma r+\frac{K}{2}r(1-r^2)$，稳态 $r_0=\sqrt{1-2\gamma/K}$，临界耦合
$\Kc=2\gamma$。而经典 Kuramoto 结果 $\Kc=2/[\pi g(\wbar)]$ 对洛伦兹
$g(\wbar)=1/(\pi\gamma)$ 给出 $\Kc=2\gamma$——\emph{一致}。

代入 \cref{eq:Hfield} 并转入旋转参考系 $w=ze^{-\ii\Om t}$（完整代数见
\cref{app:rotating}）：
%
\begin{equation}\label{eq:rot}
\keyeq{\dot w=(\ii\Delta-\gamma)w+\frac{K}{2}\bigl(w-|w|^2w\bigr)
       +\frac{F}{2}\bigl(1-w^2\bigr)},
\qquad \Delta\equiv\wbar-\Om .
\end{equation}

\noindent\textbf{无穷维系统精确降为二维 ODE。} 这不是近似，而是洛伦兹分布下的
严格结果。工程意义：二维参数扫描（$200\times200$ 网格）从数小时降到数秒，参数
拟合从不可行变为可行。\citet{Childs2008} 已给出该系统的完整分岔结构，包括
Hopf、鞍结、Takens--Bogdanov 分岔的精确位置，并指出其相图与弱非线性受迫
van der Pol 振子高度相似。

%-----------------------------------------------------------------------
\section{极坐标形式与不动点}
\label{sec:polar}

令 $w=\rho e^{\ii\Phi}$，分离 \cref{eq:rot} 的实虚部（推导见
\cref{app:polar}）：
%
\begin{subequations}\label{eq:polar}
\begin{align}
\dot\rho &= -\gamma\rho+\frac{K}{2}\rho(1-\rho^2)
            +\frac{F}{2}(1-\rho^2)\cos\Phi , \label{eq:polar-a}\\[2pt]
\dot\Phi &= \Delta-\frac{F}{2}\left(\rho+\frac{1}{\rho}\right)\sin\Phi .
            \label{eq:polar-b}
\end{align}
\end{subequations}
%
$\rho$ 是旋转系中的序参量模，即\emph{与刺激的相位锁定程度}，直接对应实验中的
试次间相干（inter-trial coherence, ITC）；$\Phi$ 是群体相对刺激的相位滞后。

\subsection{三个极限检验（代码必须硬断言）}
\label{sec:limits}

\paragraph{(i) $F=0$（无刺激）.}
$\dot\rho=\rho\bigl[\frac{K}{2}-\gamma-\frac{K}{2}\rho^2\bigr]$，故
%
\begin{equation}\label{eq:rho0}
\rho_0=\sqrt{1-\frac{2\gamma}{K}}\quad(K>\Kc),
\qquad \Kc=2\gamma .
\end{equation}

\paragraph{(ii) $K=0,\ \Delta=0$（纯受迫）.}
由 \cref{eq:polar-b} 得 $\sin\Phi=0\Rightarrow\Phi=0$；代入 \cref{eq:polar-a}：
$\frac{F}{2}\rho^2+\gamma\rho-\frac{F}{2}=0$，取正根
%
\begin{equation}\label{eq:pureforced}
\keyeq{\rho^{*}=\frac{\sqrt{\gamma^2+F^2}-\gamma}{F}}
\end{equation}
%
两端渐近：$F\ll\gamma$ 时 $\rho^{*}\to F/(2\gamma)$（\emph{线性}）；
$F\gg\gamma$ 时 $\rho^{*}\to1$（\emph{饱和}）。

\begin{remark}[第一个零参数形状预测]\label{rem:ITC}
\cref{eq:pureforced} 可直接对照真实数据。\citet{Morera2026} 测得 ITC 峰值随
刺激强度（$3.5$–$7\,$mA）单调上升，且相关强度在 TS 组与对照组间无差异；而
$3.5$–$7\,$mA 全部位于阈下／近阈（该研究报告运动阈值均值：对照
$7.68\,$mA，TS $8.03\,$mA），即处在 \cref{eq:pureforced} 的\emph{线性段}。
模型在此区间预测线性，与观测符合，且这条曲线是标定 $F(I)$ 映射的唯一入口。
\end{remark}

\paragraph{(iii) $\gamma\to0,\ K\to0$} 应退回 M0 的 Adler 方程 \cref{eq:adler}。

%-----------------------------------------------------------------------
\section{主方程：锁定态的隐式解}
\label{sec:mainEq}

锁定态即 \cref{eq:polar} 的不动点。由两式分别解出
%
\begin{equation}\label{eq:sincos}
\sin\Phi=\frac{2\Delta\rho}{F(1+\rho^2)},
\qquad
\cos\Phi=\frac{\rho\bigl[2\gamma-K(1-\rho^2)\bigr]}{F(1-\rho^2)} .
\end{equation}
%
利用 $\sin^2\Phi+\cos^2\Phi=1$ 消去 $\Phi$：
%
\begin{equation}\label{eq:main}
\keyeq{\frac{4\Delta^2\rho^2}{(1+\rho^2)^2}
      +\frac{\rho^2\bigl[2\gamma-K(1-\rho^2)\bigr]^2}{(1-\rho^2)^2}=F^2}
\end{equation}

\noindent\cref{eq:main} 是本章的\textbf{主方程}。给定 $(\Delta,K,\gamma,F)$ 解出
$\rho\in(0,1)$ 即得锁定强度。它是高次代数方程，一般需数值求根，但\emph{完全
无需模拟}。Tongue 边界由 \cref{eq:main} 与 $\partial(\cdot)/\partial\rho=0$
联立给出。

\begin{remark}[边界类型的警告]\label{rem:boundary}
\citet{Childs2008} 证明该系统的边界并非处处为鞍结分岔，参数空间中还存在
Hopf 边界、双稳区与漂移／行波态。数值扫描时\textbf{不可假定边界只有一种
类型}，需对照其完整相图核对。
\end{remark}

%-----------------------------------------------------------------------
\section{两个解析可解区制}
\label{sec:regimes}

\cref{eq:main} 虽精确但不可读。以下给出两个可解极限，\textbf{并明确各自适用
域}——混用二者会导致错误结论（见 \cref{rem:nocritdiv}）。

\subsection{区制 I：亚临界弱驱动 $\Rightarrow$ 洛伦兹调谐曲线}
\label{sec:regimeI}

条件：$K<\Kc$（静息态未自发同步，$\rho_0=0$）且 $\rho\ll1$。此时
$(1\pm\rho^2)^2\to1$、$K(1-\rho^2)\to K$，\cref{eq:main} 化为
$\rho^2\bigl[4\Delta^2+(2\gamma-K)^2\bigr]=F^2$。定义
%
\begin{equation*}
A\;\equiv\;2\gamma-K\;>\;0 ,
\end{equation*}
%
得
%
\begin{equation}\label{eq:regimeI}
\keyeq{\rho(\Delta)=\frac{F}{\sqrt{A^2+4\Delta^2}}}
\end{equation}
%
功率形式为标准洛伦兹型：
%
\begin{equation}\label{eq:power}
\rho^2(\Delta)=\frac{F^2}{A^2+4\Delta^2} .
\end{equation}
%
由此\emph{三个量同时读出}：
%
\begin{equation}\label{eq:three}
\Ostar=\wbar ,
\qquad
W=\frac{A}{2}=\gamma-\frac{K}{2} ,
\qquad
\rho_{\max}=\frac{F}{A}=\frac{F}{2\gamma-K} ,
\end{equation}
%
其中 $W$ 为\emph{功率曲线的半高半宽}（等价于 $\rho$ 的 $1/\sqrt2$ 半宽）。
归一化调谐曲线：
%
\begin{equation}\label{eq:normtune}
\frac{\rho(\Delta)}{\rho(0)}=\frac{1}{\sqrt{1+(\Delta/W)^2}} .
\end{equation}

\subsubsection{一条守恒关系}

$\gamma$ 与 $K$ 同时决定宽度与高度，且方向一致：$W\propto A$ 而
$\rho_{\max}\propto1/A$，故
%
\begin{equation}\label{eq:conserv}
\keyeq{W\cdot\rho_{\max}=\frac{F}{2}=\text{const}}
\end{equation}

\noindent\cref{eq:conserv} 表明：\textbf{调谐窗口宽度与疗效天花板成反比，二者
之积只由刺激强度决定。} 临床含义清晰而非平凡：
%
\begin{itemize}[leftmargin=2em,itemsep=2pt]
\item \textbf{窄而高的 mu 峰个体（小 $A$）}：疗效天花板高，但频率窗口窄——
      \emph{最需要个体化}。
\item \textbf{宽而平的个体（大 $A$）}：对频率不敏感，固定 $10\,$Hz 已足够，但
      天花板低。
\end{itemize}

这与实测规律逐项对应。\citet{Trajkovic2024} 报告：静息 alpha 曲线的高度预测
个体 alpha 频率（IAF）刺激下的幅度增强，而曲线越宽，移动内源频率的调制效应
越大。映射关系为\textbf{峰高 $\leftrightarrow K$，峰宽 $\leftrightarrow\gamma$}。
该经验规律此前是描述性的；\cref{eq:three,eq:conserv} 给了它一个推导。

\begin{remark}[量化示例]\label{rem:quant}
某个体 $\wbar=11\,$Hz，$W=0.5\,$Hz。固定 $10\,$Hz 时 $\Delta=1\,$Hz：
%
\begin{equation*}
\frac{\rho}{\rho_{\max}}=\frac{1}{\sqrt{1+(1/0.5)^2}}=\frac{1}{\sqrt5}=0.45 ,
\end{equation*}
%
即\emph{损失 $55\%$ 的可达带动量}。这是本课题摘要中应出现的数字量级。
\end{remark}

\paragraph{适用域声明.}
区制 I 要求 $F\ll A$（由 $\rho\ll1$ 在 $\Delta=0$ 处推出）。$K\to\Kc$ 时
$A\to0$，线性化失效，必须回到 \cref{eq:main}。

\begin{remark}[一个需要诚实交代的张力]\label{rem:tension}
$K<\Kc$ 时 mean-field 给出 $\rho_0=0$，即静息态无 mu 同步——这与“静息 mu
节律确实存在”矛盾。三种可能的和解：(a) 观测到的静息 mu 功率来自有限尺度
涨落（$\rho_0\sim N^{-1/2}$）；(b) 真实 mu 是间歇性爆发，系统在临界两侧
往复；(c) 系统实际处于超临界。\textbf{这一张力不应掩盖，它恰好指向
区制 II。}
\end{remark}

\subsection{区制 II：超临界 $\Rightarrow$ Adler 楔形与真阈值}
\label{sec:regimeII}

条件：$K>\Kc$，静息态已自发同步（$\rho_0=\sqrt{1-2\gamma/K}$ 不小），弱驱动
只作用于集体相位。此时 $\rho\approx\rho_0$ 近似固定，\cref{eq:polar-b}
退化为 Adler 方程：
%
\begin{equation}\label{eq:adlereff}
\dot\Phi=\Delta-\Feff\sin\Phi ,
\qquad
\Feff=\frac{F}{2}\left(\rho_0+\frac{1}{\rho_0}\right).
\end{equation}
%
锁定条件成为\emph{真阈值}：
%
\begin{equation}\label{eq:threshold}
\keyeq{|\Delta|\le\frac{F}{2}\left(\rho_0+\frac{1}{\rho_0}\right)}
\end{equation}

两个区制的定性差别是本章最重要的可判伪对比，汇总于 \cref{tab:regimes}。

\begin{table}[htbp]
\centering
\caption{两个区制的可判伪对比}
\label{tab:regimes}
\small
\begin{tabularx}{\textwidth}{@{}l X X@{}}
\toprule
 & 区制 I（亚临界） & 区制 II（超临界） \\
\midrule
静息态        & 未同步，$\rho_0=0$        & 自发同步，$\rho_0>0$ \\
调谐曲线      & 光滑洛伦兹，\textbf{无阈值} & 楔形，\textbf{有硬阈值} \\
剂量--反应    & 渐变                      & 窗口内有效／窗口外近乎为零 \\
个体化收益    & 渐进改善                  & \textbf{窗口外为全无效} \\
群体应答分布  & 单峰连续                  & \textbf{双峰} \\
\bottomrule
\end{tabularx}
\end{table}

\cref{eq:threshold} 给出应答者／非应答者二分的一个干净机制：若静息感觉运动
mu 处于超临界（mu 的爆发式特征部分支持如此），则 $|\wbar_i-10\,\mathrm{Hz}|$
超出该窗口的个体本质上得不到任何带动效应。由此得到一条明确的、区分两区制的
经验预测：
%
\begin{quote}
应答者与非应答者在 $|\wbar_i-10\,\mathrm{Hz}|$ 上的差异应呈\emph{阈值型}
（区制 II）而非\emph{平滑单调型}（区制 I）。
\end{quote}
%
试验实测的活跃组抽动频率降幅 $-15.6\pm36.6\,\TPM$（标准差是均值的两倍以上）
以及 $59\%$ 的应答率~\citep{Maiquez2023}——这种方差远大于均值的分布形态，与
双峰混合更相容，是支持区制 II 的间接证据。第~4 章将用虚拟队列定量检验。

\paragraph{适用域声明.}
\cref{eq:threshold} 要求 $\rho_0$ 不小。$\rho_0\to0$（即 $K\to\Kc^{+}$）时
$\Feff\to\infty$，这是 $\rho$ 固定近似失效的\emph{伪发散}。

\subsection{临界邻域：必须数值处理}
\label{sec:critical}

\begin{remark}[不存在“临界处最易带动”的简单结论]\label{rem:nocritdiv}
$|K-\Kc|\lesssim F$ 时两个近似均失效。此区域的特征是：锁相范围变宽
（\cref{eq:threshold} 中 $\Feff$ 增大），但达成的 $\rho$ 变小
（\cref{eq:regimeI} 中 $\rho_{\max}$ 随 $A\to0$ 的行为被非线性截断）——两个
效应\emph{方向相反}。因此必须直接数值求解 \cref{eq:main}。决定疗效的量
始终是 $\rho^{*}(\Om)$，“是否锁定”只是次级特征。
\end{remark}

%-----------------------------------------------------------------------
\section{无量纲化：物理只有三个参数}
\label{sec:nondim}

\begin{equation}\label{eq:nondim}
\kappa=\frac{K}{2\geff},
\qquad
f=\frac{F}{2\geff},
\qquad
\delta=\frac{\Delta}{\geff} .
\end{equation}
%
主方程 \cref{eq:main} 变为
%
\begin{equation}\label{eq:nondimmain}
\frac{\delta^2\rho^2}{(1+\rho^2)^2}
+\frac{\rho^2\bigl[1-\kappa(1-\rho^2)\bigr]^2}{(1-\rho^2)^2}=f^2 ,
\end{equation}
%
区制 I 的解变为极简形式：
%
\begin{equation}\label{eq:nondimI}
\keyeq{\rho=\frac{f}{\sqrt{(1-\kappa)^2+\delta^2}}},
\qquad
\frac{W}{\geff}=1-\kappa .
\end{equation}
%
同步相变位于 $\kappa=1$。整个模型的物理只由 $(\kappa,f,\delta)$ 三个无量纲数
决定：\textbf{内源相干度、刺激强度、频率失配}。这是本章的最简表述，也是所有
数值扫描应采用的坐标。

%-----------------------------------------------------------------------
\section{噪声的处理}
\label{sec:noise}

\textbf{必须明确交代}：OA 降维在 $D\neq0$ 时\emph{不精确}。三种处理，按严格性
递增。

\paragraph{(a) 有效展宽（工作模型）.}
\begin{equation}\label{eq:geff}
\keyeq{\geff=\gamma+D}
\end{equation}
%
依据有两条独立来源：一是 \cref{rem:Dposition} 指出的线性响应阶位置一致；二是
临界耦合的精确插值——含噪声 Kuramoto 的非相干态线性稳定性条件为
$\frac{K}{2}\Rea\int\frac{g(\omega)\,\dd\omega}{D+\ii(\omega-\wbar)}=1$，对
洛伦兹分布该积分精确等于 $1/(D+\gamma)$，故
%
\begin{equation*}
\Kc=2(\gamma+D)\qquad\text{（精确）},
\end{equation*}
%
详见 \cref{app:noise}。\textbf{适用域：$\rho\ll1$ 的线性响应区}；完全锁定态
下 \cref{eq:geff} 仅为近似。

\paragraph{(b) 圆累积量截断.}
保留到二阶累积量，得含噪声严格的四维 ODE~\citep{Tyulkina2018}，用于量化 (a)
的偏差。

\paragraph{(c) 直接 $N$ 振子 SDE.}
$N=10^3$–$10^4$，作为最终判据。

\subsection{一个零参数预测}
\label{sec:quench}

\citet{Gialopsou2025} 报告：TS 组对单脉冲正中神经刺激的试次间变异显著升高，
但在刺激作用下被\emph{归一化}。

受迫 Kuramoto 天然给出这一结果，\textbf{无需任何拟合}：外部驱动使
\cref{eq:sde} 的洗衣板势加深，相位分布收窄，$\Var(\rho)$ 必然下降。TS 的高 $D$
对应静息态高变异；施加 $F$ 后变异被“淬灭”。这是本模型第二个零参数定性预测
（第一个是 \cref{rem:ITC}），也是模型可信度的即时来源。

%-----------------------------------------------------------------------
\section{最优频率的精确位置}
\label{sec:optfreq}

由 \cref{sec:mainEq}，\cref{eq:main} 中 $\Delta$ 仅以 $\Delta^2$ 出现，故
$\rho^{*}$ 是 $\Delta$ 的偶函数：
%
\begin{equation}\label{eq:optfreq}
\keyeq{\Ostar=\wbar}
\qquad\text{（对称 } g \text{、频率平坦传输下\emph{精确}成立）}
\end{equation}
%
偏移只能来自频率依赖的传输函数 $\Xi(\Om)$。设
$\rho=F_0\Xi(\Om)/\sqrt{A^2+4(\wbar-\Om)^2}$，令
$\dd\ln\rho/\dd\Om=0$ 并对小偏移线性化（推导见 \cref{app:offset}）：
%
\begin{equation}\label{eq:offset}
\Ostar=\wbar+\frac{A^2}{4}\cdot\frac{\Xi'}{\Xi}\bigg|_{\wbar} .
\end{equation}
%
低通传输（$\Xi'<0$）使最优频率\emph{下移}，且偏移量按调谐宽度的\emph{平方}
缩放。

\begin{remark}[一个有实用价值的正面结论]\label{rem:offsetsmall}
数值代入（一阶低通、转折频率 $25\,$Hz；$W$ 取 $0.5$–$2\,$Hz）：偏移量为
$0.003$–$0.06\,$Hz，\emph{远小于} EEG 峰频估计精度（典型 $\pm0.2$–$0.5\,$Hz）。
因此开环生物标记算法可直接取 $\Ostar_i=\hat{\wbar}_i$，\textbf{无需模型修正
项}。第~3 章的算法设计以此为依据。
\end{remark}

%-----------------------------------------------------------------------
\section{脉冲驱动与高阶锁定结构}
\label{sec:pulse}

设备输出为脉冲串而非正弦波。理想脉冲串的傅里叶展开（$\Om=2\pi/T$）：
%
\begin{equation}\label{eq:pulse}
F_p\sum_{n}\delta(t-nT)
=\frac{F_p}{T}\left[1+2\sum_{m=1}^{\infty}\cos(m\Om t)\right].
\end{equation}

\paragraph{次谐波驱动 tongue.}
第 $m$ 次谐波在 $m\Om\approx\wbar$ 时可 $1{:}1$ 锁定，即在 $\Om\approx\wbar/m$
处出现附加 tongue。例如 $5\,$Hz 脉冲串的二次谐波可锁定 $10\,$Hz mu。这
\emph{不需要}超出 \cref{sec:reduction} 三步约化的任何假设，是脉冲驱动的直接
后果。

\paragraph{高阶 $p{:}q$ tongue.}
需要非正弦相位耦合函数或振幅动力学（Stuart--Landau），超出本章约化范围，
列为第~6 章展望。

\subsection{一个判据及其局限}
\label{sec:subharm}

线性时不变系统受 $\Om$ 驱动只能输出 $m\Om$，\textbf{无法输出 $\Om/2$}。因此
真次谐波响应是非线性锁定的\emph{充分}判据，可原则上区分“真带动”与“诱发反应
线性叠加”——后者正是 \citet{Houlgreave2022} 对 $12\,$Hz 条件的结论：
$20\,$Hz 刺激产生了 beta 波段的真正带动，而 $12\,$Hz 的响应主要归因于节律性
感觉诱发电位。

\citet{Morera2026} 的时频分析显示 $10\,$Hz 刺激下 $5\,$Hz 与 $10\,$Hz 功率
同时上升，且上升与 MNS 起始同时发生。

\begin{remark}[对该观测的诚实评估：先验解释为起始瞬变]\label{rem:onset}
该实验采用 $1\,$秒脉冲串 $+\,5\,$秒间隔；$1\,$秒矩形窗的频谱宽度约 $1\,$Hz，
而\emph{串起始的宽带瞬变在连续小波变换中必然表现为低频功率上升，包括
$5\,$Hz 附近}。原文“与起始同时发生”的措辞本身即指向起始瞬变。判据在理论上
有效，但该具体观测的先验解释是 onset transient。若要真正检验，需要：
(i) 剔除起始 $300\,$ms 后分析稳态段；(ii) 采用连续长时程刺激而非短串；
(iii) 检验相位一致性而非仅功率。\textbf{本条列为附录中的候选实验设计
（\cref{app:subharm}），不作为主预测。}
\end{remark}

%-----------------------------------------------------------------------
\section{观测模型：从 $\rho$ 到可测量}
\label{sec:observation}

模型的经济性在于三条观测链路共用同一个 $\rho$：

\begin{table}[htbp]
\centering
\caption{三条观测链路}
\label{tab:obs}
\small
\begin{tabularx}{\textwidth}{@{}l l X@{}}
\toprule
观测量 & 模型对应 & 数据来源 \\
\midrule
试次间相干 ITC & $\rho$ & \citet{Morera2026}; \citet{Houlgreave2022} \\
诱发 mu 功率   & $\propto\rho^2$ & \citet{Houlgreave2022} \\
抽动频率 \TPM  & \cref{eq:lambda} & \citet{Maiquez2023}; \citet{Iverson2023a} \\
\bottomrule
\end{tabularx}
\end{table}

\subsection{机制推导}
\label{sec:mechanism}

设运动“就绪度” $m(t)$ 为零均值、方差 $\sigma_m^2$ 的平稳高斯过程（纹状体噪声
驱动，对应 TS 中抑制性中间神经元减少导致 CSTC 环路噪声增高的病理生理
假设~\citep{Albin2006}）。$m(t)$ 越过阈值 $\Theta$ 即触发抽动。由 Rice 越阈
公式，高阈值下越阈率
%
\begin{equation}\label{eq:rice}
\lambda=\nu_0\exp\!\left(-\frac{\Theta^2}{2\sigma_m^2}\right),
\qquad
\nu_0=\frac{1}{2\pi}\frac{\sigma_{\dot m}}{\sigma_m} .
\end{equation}
%
\cref{hyp:H2} 的具体形式为：mu 相干抬高有效阈值，
%
\begin{equation}\label{eq:thetaeff}
\Thf=\Theta_0+a\rho .
\end{equation}
%
代入 \cref{eq:rice}（展开见 \cref{app:rice}）并\emph{保留 $\rho$ 的一阶项}，
即得唯象形式，且系数有了明确出处：
%
\begin{equation}\label{eq:lambda}
\keyeq{\lambda=\lambda_0e^{-\beta\rho}},
\qquad
\beta=\frac{a\Theta_0}{\sigma_m^2} .
\end{equation}

\begin{remark}[两条链路的统一]\label{rem:unify}
\cref{eq:lambda} 的唯象形式与 \cref{eq:rice}--\cref{eq:thetaeff} 的机制形式
\emph{不是并列的两个选项，而是同一推导的一阶截断与完整形式}。被丢弃的
$\rho^2$ 项意味着真实依赖略\emph{陡于}指数——这是一条可检验的二阶偏差。

同时注意 $\beta\propto1/\sigma_m^2$：\emph{噪声越大，门控效率越低}。TS 的高
噪声一方面通过 $D$ 压低 $\rho$，另一方面通过 $\sigma_m$ 压低 $\beta$，
\textbf{两条路径同向削弱疗效}。这是模型对“为何 TS 反应异质性大”的一个结构性
解释。
\end{remark}

\subsection{疗效函数}
\label{sec:efficacy}

\begin{equation}\label{eq:efficacy}
E(\Om)=1-\frac{\langle\lambda\rangle_{\mathrm{stim}}}
                {\langle\lambda\rangle_{\mathrm{base}}}
      =1-\exp\bigl\{-\beta\bigl[\rho^{*}(\Om)-\rho_0\bigr]\bigr\} .
\end{equation}
%
区制 I（$\rho_0=0$）下代入 \cref{eq:regimeI}，得\textbf{完整的个体调谐曲线
闭式解}：
%
\begin{equation}\label{eq:Ei}
\keyeq{E_i(\Om)=1-\exp\!\left[-\frac{\beta_iF_i}
       {\sqrt{A_i^2+4(\wbar_i-\Om)^2}}\right]}
\end{equation}
%
\cref{eq:Ei} 是第~3 章个体化算法的目标函数，也是第~4 章虚拟队列的前向模型。
作为 Cox 点过程的条件强度，它可直接对真实 \TPM{} 计数数据做极大似然估计。

\subsection{一条衍生预测}
\label{sec:phasegate}

若 \cref{eq:thetaeff} 中保留相位依赖，
$\Thf=\Theta_0+a\rho+b\rho\cos\varphi_{\mathrm{gate}}(t)$，则
%
\begin{quote}
\textbf{预测}：抽动发生概率在 mu 周期内呈相位调制。
\end{quote}
%
这可用带时间戳的抽动记录直接检验，并构成“相位锁定闭环刺激”（比调频率更进
一步）的理论依据。

%-----------------------------------------------------------------------
\section{慢变量层：online 与 offline 必须分离}
\label{sec:slow}

一个不可回避的逻辑区分：
%
\begin{itemize}[leftmargin=2em,itemsep=2pt]
\item \textbf{Online 效应}（刺激期间 \TPM{} 下降）：由 $\rho(t)$ 直接给出，
      即时、不外溢。与 \citet{Iverson2023a} 报告 MNS 益处在刺激结束后不
      持续一致。
\item \textbf{Offline 效应}（4 周后 \YGTSS{} 降 $7.13$ 分）：\emph{带动本身
      无法解释}，因为 $\rho$ 在刺激停止后立即回落。必须引入慢变量。
\end{itemize}

最简 Hebbian 形式（时标 $\tau_p\sim$ 天）：
%
\begin{equation}\label{eq:hebb}
\tau_p\frac{\dd K}{\dd t}=-(K-K_0)+\eta\,\overline{\rho}
\qquad\text{或}\qquad
\tau_p\frac{\dd D}{\dd t}=-(D-D_0)-\eta'\overline{\rho} ,
\end{equation}
%
其中 $\overline{\rho}$ 为单次 session 的时间平均。累积刺激缓慢提升 $\kappa$
（或降低 $D$），进而提升基线 $\rho_0$，降低基线抽动率。

\subsection{统一两组看似矛盾的实验结果}
\label{sec:unifyexp}

因 offline 效应对 $\overline{\rho}$ 做\emph{时间积分}，而 $\overline{\rho}$
本身按 \cref{eq:normtune} 随 $|\Delta|$ 衰减，积分放大了去谐惩罚：
%
\begin{quote}
\textbf{预测}：offline 效应的频率选择性\emph{强于} online 效应。
\end{quote}
%
这解释了为何 \citet{Iverson2023a} 的单次交叉实验看不到节律／非节律差异
（online，选择性弱），而 \citet{Maiquez2023} 的 4 周方案能与假刺激显著分离
（offline，选择性强）。\textbf{模型给出了这两组结果的统一解释}，这是第~5 章
的独立结果小节。

\subsection{适应项}
\label{sec:adapt}

\citet{Morera2026} 显示 ITC 增幅在刺激期内衰减，且 TS 组的衰减大于对照组。
建模为
%
\begin{equation}\label{eq:adapt}
F(t)=F_0\bigl[a_\infty+(1-a_\infty)e^{-t/T_a}\bigr],
\end{equation}
%
TS 组 $T_a$ 更短。这是可直接拟合的两参数动态特征，同时为设备采用
$2\,$min on $/\,1\,$min off 占空比提供了理论依据（允许 $F$ 在间歇期恢复）。

%-----------------------------------------------------------------------
\section{参数与可辨识性}
\label{sec:identifiability}

表面上每个体有 9 个参数
$\{\wbar,\gamma,K,D,F,\beta,\lambda_0,T_a,a_\infty\}$。但在区制 I 下，由
\cref{eq:Ei} 可见\textbf{只有三个可辨识组合}：
%
\begin{equation}\label{eq:ident}
\keyeq{\Theta_i=\bigl\{\;\wbar_i,\quad
       A_i\equiv2\gamma_{\mathrm{eff},i}-K_i,\quad
       B_i\equiv\beta_iF_i\;\bigr\}}
\end{equation}
%
因为 $K$ 与 $\geff$ 只以差值 $A$ 出现，$F$ 与 $\beta$ 只以乘积 $B$ 出现
（Fisher 信息矩阵的零方向见 \cref{app:ident}）。

\begin{remark}[决定项目可行性的结果]\label{rem:identgood}
这一结果\emph{消除了}对 $K$--$F$ 共线性的担忧：不是“报告参数不可辨识”，而是
模型在该区制下本质上只有三个自由度。这使个体拟合从 9 参数问题降为 3 参数
问题——在每人仅有数个 session 数据的临床现实下，这是决定项目可行性的关键。
\end{remark}

三个参数的经验先验：

\begin{table}[htbp]
\centering
\caption{三个可辨识参数的经验先验来源}
\label{tab:priors}
\small
\begin{tabularx}{\textwidth}{@{}l X@{}}
\toprule
参数 & 先验来源 \\
\midrule
$\wbar_i$ & 峰值 alpha 频率（PAF）发育轨迹：13 岁约 $9.7\,$Hz，渐近
            $10.1\,$Hz~\citep{Cellier2022}；成人感觉运动 mu 去同步峰
            $10$–$12\,$Hz~\citep{Berchicci2011}；个体标准差约 $1\,$Hz \\
$A_i$     & specparam 提取的 mu 峰宽（典型 FWHM $1.5$–$3\,$Hz）；TS 组 $D$ 更高 \\
$B_i$     & 由 \cref{eq:pureforced} 标定 $F(I)$；$\beta$ 由 \TPM{} 剂量反应拟合 \\
\bottomrule
\end{tabularx}
\end{table}

\begin{remark}[方法学要求，非可选]\label{rem:fooof}
$\wbar_i$ 必须用 specparam/FOOOF~\citep{Donoghue2020} 剥离 $1/f$ 非周期成分后
再取峰。\citet{Adelhofer2021} 已证明 TS 在感觉运动整合中的无标度／非周期活动
升高，直接取 FFT 峰会引入\emph{与组别相关}的系统偏差——而该偏差恰好污染
本模型最核心的参数。
\end{remark}

%-----------------------------------------------------------------------
\section{可判伪预测清单}
\label{sec:predictions}

\begin{table}[htbp]
\centering
\caption{十条可判伪预测及其当前证据状态}
\label{tab:predictions}
\footnotesize
\begin{tabularx}{\textwidth}{@{}l X l X l@{}}
\toprule
\# & 预测 & 依据 & 检验途径 & 状态 \\
\midrule
P1 & 最优频率 $=$ 个体 mu 峰频，偏移 $<0.1\,$Hz
   & \cref{eq:optfreq,eq:offset} & 闭环调频实验 & 待检验 \\
P2 & 调谐宽度随强度展开 $\Rightarrow$ 低强度方案最需个体化
   & tongue 几何 & $2\times2$（强度$\times$调频） & 待检验 \\
P3 & 窄高峰个体：高天花板 $+$ 强频率敏感；$W\rho_{\max}$ 守恒
   & \cref{eq:three,eq:conserv} & 静息谱形 $\to$ 疗效相关 & \textbf{部分支持} \\
P4 & 年龄 $\to$ mu 频率 $\to$ 去谐 $\to$ 疗效
   & \cref{eq:normtune} & 试验已发表年龄数据 & 待检验 \\
P5 & 若超临界，应答／非应答呈\emph{阈值型}而非平滑型
   & \cref{eq:threshold} & $|\wbar_i-10|$ vs 应答 & \textbf{间接支持} \\
P6 & ITC--强度曲线低强度段线性、高强度饱和
   & \cref{eq:pureforced} & --- & \textbf{已支持（零参数）} \\
P7 & 刺激淬灭试次间变异
   & \cref{eq:sde} 势加深 & --- & \textbf{已支持（零参数）} \\
P8 & offline 频率选择性 $>$ online
   & \cref{eq:hebb} 积分 & 统一两组实验 & 待检验 \\
P9 & 抽动概率呈 mu 相位调制
   & \cref{eq:thetaeff} 相位项 & 抽动时间戳 & 待检验 \\
P10 & 真次谐波响应可判别带动 vs 诱发叠加
   & \cref{eq:pulse} & 需专门实验设计 & \textbf{降级，附录} \\
\bottomrule
\end{tabularx}
\end{table}

P6、P7 已被现有数据支持且\emph{无自由参数}——这是模型的即时可信度锚点。
P2、P3、P5、P8 是最具原创性的四条。

%-----------------------------------------------------------------------
\section{模型边界：明确列出它解释不了什么}
\label{sec:boundary}

学术诚实是本章能否被接受的分水岭。以下每一条都应写入正文而非藏在讨论末尾。

\begin{enumerate}[leftmargin=2.2em,itemsep=4pt]

\item \textbf{非节律 $\approx$ 节律（最强反证）.}
模型中非节律驱动对应宽带随机输入，$\rho$ 提升应远小于节律驱动。
\citet{Iverson2023a} 报告节律与非节律 $12\,$Hz MNS 对抽动频率、强度、冲动的
改善无显著差异；\citet{Houlgreave2025} 进一步发现节律与非节律刺激的紧张性
神经调质效应（GABA 下降）亦无显著差异。\emph{本模型无法解释两者等效。}
应明确写为带动假设的边界，剩余部分需归因于其他机制（同步放电占用皮层资源以
降低感觉运动噪声；或 GABA 介导的非频率特异性效应）。

\item \textbf{相位约化丢弃振幅动力学.}
真实 mu 是爆发式而非持续极限环。burst 统计、间歇性，以及 \cref{rem:tension}
中 $\rho_0=0$ 的张力，都需要升级到 Stuart--Landau 或含亚临界 Hopf 的模型。

\item \textbf{洛伦兹分布是数学便利，非生理事实.}
真实 mu 谱峰有界（$\omega>0$）、与 theta/beta 重叠、且不对称。检验方式：用
高斯或偏态 $g(\omega)$ 做 $N$ 振子模拟，量化 \cref{eq:regimeI} 的偏差。预期
偏差主要影响 tongue 尾部而非峰位（因 \cref{eq:optfreq} 只依赖对称性）。

\item \textbf{单群 mean-field 无空间结构.}
不含 CSTC 环路拓扑、不含皮层--基底神经节传导延迟环。P10 若成立可能需双群
模型才能定量。

\item \textbf{安慰剂无法排除.}
\citet{Iverson2023a} 明确指出其数据不能排除安慰剂效应。模型能提供的是
\emph{安慰剂无法模仿的频率依赖模式}（P1--P3、P5）——这正是本模型的临床
价值：它把“是否有效”的问题转化为“疗效是否按预测的方式随频率变化”，后者
对安慰剂免疫。

\item \textbf{外周抖动可忽略（负面结果）.}
$\Xi_{\mathrm{jitter}}=e^{-\sigma_\tau^2\Om^2/2}$，代入 $\sigma_\tau\le3\,$ms、
$\Om=2\pi\cdot10\,$rad/s：衰减约 $1.8\%$。需 $\sigma_\tau\gtrsim10.6\,$ms 才
产生 $20\%$ 效应，外周传入通路不现实（推导见 \cref{app:jitter}）。
\emph{结论：低频段外周路径是“干净”的传输通道，该机制仅在 gamma 波段方案中
需考虑。}

\end{enumerate}

%-----------------------------------------------------------------------
\section*{本章小结}
\addcontentsline{toc}{section}{本章小结}
\label{sec:summary}

从三步约化出发，经 OA 精确降维，得到主方程 \cref{eq:main}。两个可解区制给出
核心结果：
%
\begin{align*}
\text{区制 I：}\quad &\rho=\frac{f}{\sqrt{(1-\kappa)^2+\delta^2}},
  \qquad W\cdot\rho_{\max}=\frac{F}{2};\\
\text{区制 II：}\quad &|\Delta|\le\frac{F}{2}\left(\rho_0+\frac{1}{\rho_0}\right).
\end{align*}
%
经观测模型 \cref{eq:lambda} 转为疗效函数 \cref{eq:Ei}，参数压缩为三个可辨识量
\cref{eq:ident}，产出十条可判伪预测——其中两条已被现有数据零参数支持。

模型的核心论断可以用一句话概括：

\begin{quote}
\textbf{该试验已在 Arnold tongue 的振幅轴上实施了个体化（按运动阈值设定
$120\%$ 强度），却在频率轴上采用固定值；而由 \cref{eq:three} 与
\cref{eq:conserv}，频率轴恰是锁定条件中更敏感的那一维。}
\end{quote}


\appendix
%=======================================================================
%  附录 A  完整代数推导
%=======================================================================
\chapter{受迫 Kuramoto 模型的完整推导}
\label{app:derivation}

本附录给出第~\ref{chap:model} 章中所有非平凡代数步骤的完整推导，以便读者独立
复核。所有结果均可用符号计算软件逐行验证；配套的验证脚本见开源仓库
\texttt{tests/test\_analytic.py}。

%-----------------------------------------------------------------------
\section{从 OA 方程到旋转参考系}
\label{app:rotating}

起点为 Ott--Antonsen 方程 \cref{eq:oa}，
%
\begin{equation}
\dot z=(\ii\wbar-\gamma)z+\tfrac12\bigl(H-H^{*}z^{2}\bigr),
\end{equation}
%
其中等效驱动场为 \cref{eq:Hfield}：$H=Kz+Fe^{\ii\Om t}$，故
$H^{*}=Kz^{*}+Fe^{-\ii\Om t}$。代入：
%
\begin{align}
\dot z
&=(\ii\wbar-\gamma)z
 +\frac12\Bigl[Kz+Fe^{\ii\Om t}-\bigl(Kz^{*}+Fe^{-\ii\Om t}\bigr)z^{2}\Bigr]
 \notag\\
&=(\ii\wbar-\gamma)z
 +\frac{K}{2}\bigl(z-|z|^{2}z\bigr)
 +\frac{F}{2}\bigl(e^{\ii\Om t}-e^{-\ii\Om t}z^{2}\bigr),
 \label{eq:app-zdot}
\end{align}
%
其中用到 $z^{*}z^{2}=|z|^{2}z$。

现引入旋转系变量
%
\begin{equation}
w=z\,e^{-\ii\Om t}
\quad\Longleftrightarrow\quad
z=w\,e^{\ii\Om t},
\qquad
\dot z=\bigl(\dot w+\ii\Om w\bigr)e^{\ii\Om t}.
\end{equation}
%
逐项代入 \cref{eq:app-zdot}，注意 $|z|^{2}=|w|^{2}$（模不变）以及
%
\begin{equation}
e^{-\ii\Om t}z^{2}=e^{-\ii\Om t}\cdot w^{2}e^{2\ii\Om t}=w^{2}e^{\ii\Om t},
\end{equation}
%
得
%
\begin{equation}
\bigl(\dot w+\ii\Om w\bigr)e^{\ii\Om t}
=(\ii\wbar-\gamma)w\,e^{\ii\Om t}
+\frac{K}{2}\bigl(w-|w|^{2}w\bigr)e^{\ii\Om t}
+\frac{F}{2}\bigl(1-w^{2}\bigr)e^{\ii\Om t}.
\end{equation}
%
两边除以 $e^{\ii\Om t}$ 并把 $\ii\Om w$ 移到右侧：
%
\begin{equation}
\dot w=\bigl[\ii(\wbar-\Om)-\gamma\bigr]w
+\frac{K}{2}\bigl(w-|w|^{2}w\bigr)
+\frac{F}{2}\bigl(1-w^{2}\bigr),
\end{equation}
%
即 \cref{eq:rot}，其中 $\Delta=\wbar-\Om$。\hfill$\blacksquare$

\begin{remark}
关键点在于旋转变换把\emph{显式时间依赖}完全消去：驱动项 $Fe^{\ii\Om t}$ 变为
常数 $F$，$e^{-\ii\Om t}z^2$ 项变为 $w^2$。这使系统成为\emph{自治}二维系统，
从而其全部长时行为可由不动点与极限环分类穷尽——这是 \citet{Childs2008} 能够
给出完整分岔图的前提。
\end{remark}

%-----------------------------------------------------------------------
\section{极坐标分解}
\label{app:polar}

令 $w=\rho e^{\ii\Phi}$（$\rho\ge0$），则
%
\begin{equation}
\dot w=\bigl(\dot\rho+\ii\rho\dot\Phi\bigr)e^{\ii\Phi}.
\end{equation}
%
将 \cref{eq:rot} 两边乘以 $e^{-\ii\Phi}$：
%
\begin{equation}
\dot\rho+\ii\rho\dot\Phi
=(\ii\Delta-\gamma)\rho
+\frac{K}{2}\bigl(\rho-\rho^{3}\bigr)
+\frac{F}{2}\bigl(e^{-\ii\Phi}-\rho^{2}e^{\ii\Phi}\bigr).
\end{equation}
%
展开最后一项：
%
\begin{align}
e^{-\ii\Phi}-\rho^{2}e^{\ii\Phi}
&=(\cos\Phi-\ii\sin\Phi)-\rho^{2}(\cos\Phi+\ii\sin\Phi)\notag\\
&=(1-\rho^{2})\cos\Phi-\ii\,(1+\rho^{2})\sin\Phi .
\label{eq:app-expand}
\end{align}
%
分离实虚部：

\paragraph{实部.}
\begin{equation}
\dot\rho=-\gamma\rho+\frac{K}{2}\rho(1-\rho^{2})
         +\frac{F}{2}(1-\rho^{2})\cos\Phi ,
\end{equation}
即 \cref{eq:polar-a}。

\paragraph{虚部.}
\begin{equation}
\rho\dot\Phi=\Delta\rho-\frac{F}{2}(1+\rho^{2})\sin\Phi
\quad\Longrightarrow\quad
\dot\Phi=\Delta-\frac{F}{2}\frac{1+\rho^{2}}{\rho}\sin\Phi
        =\Delta-\frac{F}{2}\left(\rho+\frac{1}{\rho}\right)\sin\Phi ,
\end{equation}
即 \cref{eq:polar-b}。\hfill$\blacksquare$

\begin{remark}
注意 \cref{eq:app-expand} 中 $(1-\rho^2)$ 与 $(1+\rho^2)$ 的\emph{符号不对称}：
振幅方程含 $(1-\rho^2)$，在 $\rho\to1$ 时驱动对振幅失效（已完全同步，无可再
增）；相位方程含 $(1+\rho^2)$，恒不为零。这一结构差异是 \cref{eq:main} 中两项
分母不同的根源，也是区制 II 中 $\Feff$ 出现 $\rho_0+1/\rho_0$ 组合的来源。
\end{remark}

%-----------------------------------------------------------------------
\section{主方程的导出}
\label{app:main}

锁定态为 \cref{eq:polar} 的不动点，即 $\dot\rho=\dot\Phi=0$。

由 $\dot\Phi=0$：
%
\begin{equation}
\Delta=\frac{F}{2}\frac{1+\rho^{2}}{\rho}\sin\Phi
\quad\Longrightarrow\quad
\sin\Phi=\frac{2\Delta\rho}{F(1+\rho^{2})} .
\label{eq:app-sin}
\end{equation}

由 $\dot\rho=0$：
%
\begin{equation}
\frac{F}{2}(1-\rho^{2})\cos\Phi=\gamma\rho-\frac{K}{2}\rho(1-\rho^{2})
=\frac{\rho}{2}\bigl[2\gamma-K(1-\rho^{2})\bigr],
\end{equation}
%
故
%
\begin{equation}
\cos\Phi=\frac{\rho\bigl[2\gamma-K(1-\rho^{2})\bigr]}{F(1-\rho^{2})} .
\label{eq:app-cos}
\end{equation}
%
\cref{eq:app-sin,eq:app-cos} 合为 \cref{eq:sincos}。代入
$\sin^{2}\Phi+\cos^{2}\Phi=1$：
%
\begin{equation}
\frac{4\Delta^{2}\rho^{2}}{F^{2}(1+\rho^{2})^{2}}
+\frac{\rho^{2}\bigl[2\gamma-K(1-\rho^{2})\bigr]^{2}}{F^{2}(1-\rho^{2})^{2}}=1 ,
\end{equation}
%
两边乘 $F^{2}$ 即得主方程 \cref{eq:main}。\hfill$\blacksquare$

\paragraph{关于解的存在性.}
\cref{eq:main} 左端在 $\rho\to1^{-}$ 时因 $(1-\rho^2)^{-2}$ 发散（除非
$2\gamma-K(1-\rho^2)\to0$ 同时发生），在 $\rho\to0^{+}$ 时趋于 $0$。故对任意
$F>0$ 至少存在一个根；根的个数与稳定性需结合 \cref{rem:boundary} 所述的完整
分岔结构判断，不可仅由单调性推断。

%-----------------------------------------------------------------------
\section{区制 I 的线性化与调谐曲线宽度}
\label{app:regimeI}

\paragraph{线性化.}
设 $\rho\ll1$。则 $(1+\rho^{2})^{2}=1+O(\rho^{2})$、
$(1-\rho^{2})^{2}=1+O(\rho^{2})$、$K(1-\rho^{2})=K+O(\rho^{2})$。保留最低阶，
\cref{eq:main} 化为
%
\begin{equation}
\rho^{2}\bigl[4\Delta^{2}+(2\gamma-K)^{2}\bigr]=F^{2}
\quad\Longrightarrow\quad
\rho=\frac{F}{\sqrt{A^{2}+4\Delta^{2}}},
\qquad A\equiv2\gamma-K .
\end{equation}

\paragraph{自洽条件.}
要求 $\rho\ll1$，在最不利的 $\Delta=0$ 处给出
%
\begin{equation}
\keyeq{F\ll A}
\end{equation}
%
这是区制 I 的定量适用域。注意 $A>0$ 等价于 $K<2\gamma=\Kc$，即亚临界条件。

\paragraph{半高半宽.}
功率 $\rho^{2}=F^{2}/(A^{2}+4\Delta^{2})$ 的峰值为 $F^{2}/A^{2}$（在
$\Delta=0$）。半高处
%
\begin{equation}
\frac{F^{2}}{A^{2}+4\Delta^{2}}=\frac{1}{2}\frac{F^{2}}{A^{2}}
\;\Longrightarrow\;
A^{2}+4\Delta^{2}=2A^{2}
\;\Longrightarrow\;
|\Delta|=\frac{A}{2}\equiv W .
\end{equation}
%
在 $|\Delta|=W$ 处，$\rho=F/\sqrt{A^{2}+A^{2}}=\rho_{\max}/\sqrt{2}$，故 $W$
同时是 $\rho$ 的 $1/\sqrt2$ 半宽——与正文 \cref{eq:three} 的表述一致。

\paragraph{守恒关系.}
\begin{equation}
W\cdot\rho_{\max}=\frac{A}{2}\cdot\frac{F}{A}=\frac{F}{2},
\end{equation}
%
即 \cref{eq:conserv}。该关系独立于 $\gamma$ 与 $K$ 的具体取值，仅依赖二者
通过 $A$ 的组合——这正是 \cref{eq:ident} 中 $A$ 为可辨识量的直接体现。
\hfill$\blacksquare$

%-----------------------------------------------------------------------
\section{区制 II 的 Adler 约化}
\label{app:regimeII}

设 $K>\Kc$ 且驱动弱，使 $\rho$ 在整个刺激过程中近似保持其自发值
$\rho_{0}=\sqrt{1-2\gamma/K}$。此时 \cref{eq:polar-a} 近似满足
（$\dot\rho\approx0$ 由 $\rho_0$ 的定义保证，$F$ 的修正为高阶），动力学退化为
\cref{eq:polar-b} 的单变量方程：
%
\begin{equation}
\dot\Phi=\Delta-\Feff\sin\Phi ,
\qquad
\Feff=\frac{F}{2}\left(\rho_{0}+\frac{1}{\rho_{0}}\right).
\end{equation}
%
这与 \cref{eq:adler} 形式相同，故锁定条件直接由 \cref{eq:lock} 给出：
%
\begin{equation}
|\Delta|\le\Feff=\frac{F}{2}\left(\rho_{0}+\frac{1}{\rho_{0}}\right),
\end{equation}
%
即 \cref{eq:threshold}。\hfill$\blacksquare$

\paragraph{两点注意.}
\begin{enumerate}[leftmargin=2.2em,itemsep=2pt]
\item 由 AM--GM 不等式，$\rho_{0}+1/\rho_{0}\ge2$，等号在 $\rho_{0}=1$ 取到。
      故 $\Feff\ge F$：\emph{自发同步总是放大集体对外部驱动的相位敏感度}。
      直观解释：同步群体作为单一宏观振子响应驱动，其有效耦合大于单个振子。
\item $\rho_{0}\to0^{+}$（即 $K\to\Kc^{+}$）时 $\Feff\to\infty$。这\textbf{不
      是}物理发散，而是 $\rho\approx\rho_0$ 固定近似失效的信号——$\rho_0$
      本身趋于零时，驱动对振幅的相对影响不再可忽略。此区域必须回到
      \cref{eq:main} 数值求解（\cref{rem:nocritdiv}）。
\end{enumerate}

%-----------------------------------------------------------------------
\section{噪声：临界耦合的精确插值}
\label{app:noise}

考虑含噪声、含频率异质性的 Kuramoto 模型 \cref{eq:micro}。非相干态
($\rho=0$) 的线性稳定性由如下特征方程决定（标准 Fokker--Planck 线性化
结果，见 \citealp{Strogatz2000}）：
%
\begin{equation}
\frac{K}{2}\int_{-\infty}^{\infty}
\frac{g(\omega)\,\dd\omega}{D+\ii(\omega-\wbar)}=1 .
\label{eq:app-stab}
\end{equation}

对洛伦兹分布 \cref{eq:lorentz}，令 $u=\omega-\wbar$：
%
\begin{equation}
\int_{-\infty}^{\infty}\frac{1}{D+\ii u}\cdot
\frac{\gamma/\pi}{u^{2}+\gamma^{2}}\,\dd u .
\end{equation}
%
被积函数在上半平面有单极点 $u=\ii\gamma$（来自洛伦兹）；$1/(D+\ii u)$ 的极点
位于 $u=\ii D$，也在上半平面。但注意
$\frac{1}{D+\ii u}=\frac{1}{\ii}\cdot\frac{1}{u-\ii D}$，故在上半平面沿大
半圆闭合，由留数定理
%
\begin{align}
\oint &= 2\pi\ii\left[
\underset{u=\ii\gamma}{\mathrm{Res}}+\underset{u=\ii D}{\mathrm{Res}}\right].
\end{align}
%
直接计算（或更简便地，利用洛伦兹分布的特征函数
$\int g(\omega)e^{\ii s u}\dd u=e^{-\gamma|s|}$ 与
$\frac{1}{D+\ii u}=\int_{0}^{\infty}e^{-(D+\ii u)s}\dd s$ 交换积分次序）：
%
\begin{equation}
\int_{0}^{\infty}e^{-Ds}\left[\int g\,e^{-\ii u s}\dd u\right]\dd s
=\int_{0}^{\infty}e^{-Ds}e^{-\gamma s}\,\dd s
=\frac{1}{D+\gamma} .
\end{equation}
%
代回 \cref{eq:app-stab}：
%
\begin{equation}
\keyeq{\Kc=2(\gamma+D)}
\end{equation}
%
\hfill$\blacksquare$

\begin{remark}[\cref{eq:geff} 的地位]\label{rem:geffstatus}
上式表明 $\geff=\gamma+D$ 在\emph{非相干态线性稳定性层面是精确的}——这比
\cref{rem:Dposition} 的启发式论证强得多。但它\textbf{不}保证在完全锁定态
（$\rho$ 接近 $1$）下仍精确：那里噪声的作用不再是简单展宽，因为
Fokker--Planck 方程的非线性项开始主导。故正文声明的适用域“$\rho\ll1$ 的线性
响应区”是必要的。定量偏差应用 \citet{Tyulkina2018} 的圆累积量方法或直接
SDE 模拟评估。
\end{remark}

%-----------------------------------------------------------------------
\section{最优频率偏移的推导与数值估计}
\label{app:offset}

设传输函数 $\Xi(\Om)$ 描述从刺激器设定频率到皮层有效驱动强度的频率依赖
衰减，则
%
\begin{equation}
\rho(\Om)=\frac{F_{0}\,\Xi(\Om)}{\sqrt{A^{2}+4(\wbar-\Om)^{2}}} .
\end{equation}
%
取对数：
%
\begin{equation}
\ln\rho=\ln F_{0}+\ln\Xi(\Om)
-\tfrac12\ln\!\bigl[A^{2}+4(\wbar-\Om)^{2}\bigr].
\end{equation}
%
求导（注意 $\dd(\wbar-\Om)/\dd\Om=-1$）：
%
\begin{equation}
\frac{\dd\ln\rho}{\dd\Om}
=\frac{\Xi'}{\Xi}
-\frac12\cdot\frac{-8(\wbar-\Om)}{A^{2}+4(\wbar-\Om)^{2}}
=\frac{\Xi'}{\Xi}+\frac{4(\wbar-\Om)}{A^{2}+4(\wbar-\Om)^{2}} .
\end{equation}
%
令 $u=\Om-\wbar$（故 $\wbar-\Om=-u$）并置零：
%
\begin{equation}
\frac{\Xi'}{\Xi}-\frac{4u}{A^{2}+4u^{2}}=0 .
\end{equation}
%
对小偏移 $|u|\ll A/2$ 线性化（分母取 $A^{2}$）：
%
\begin{equation}
u\simeq\frac{A^{2}}{4}\cdot\frac{\Xi'}{\Xi}\bigg|_{\wbar},
\end{equation}
%
即 \cref{eq:offset}。\hfill$\blacksquare$

\paragraph{数值估计.}
取一阶低通 $\Xi(f)=\bigl[1+(f/f_{c})^{2}\bigr]^{-1/2}$，转折频率
$f_{c}=25\,$Hz。则
%
\begin{equation}
\frac{\Xi'}{\Xi}=-\frac{f/f_{c}^{2}}{1+(f/f_{c})^{2}}
\;\xrightarrow[\;f=10\,\mathrm{Hz}\;]{}\;
-\frac{10/625}{1+0.16}=-0.0138\ \mathrm{Hz^{-1}} .
\end{equation}
%
代入 \cref{eq:offset}，见 \cref{tab:offset}。

\begin{table}[htbp]
\centering
\caption{最优频率偏移的数值估计（一阶低通，$f_c=25\,$Hz）}
\label{tab:offset}
\small
\begin{tabular}{@{}lccc@{}}
\toprule
$W$ (Hz) & $A=2W$ (Hz) & $A^{2}/4$ & 偏移 $|u|$ (Hz) \\
\midrule
$0.5$ & $1.0$ & $0.25$ & $0.003$ \\
$1.0$ & $2.0$ & $1.00$ & $0.014$ \\
$2.0$ & $4.0$ & $4.00$ & $0.055$ \\
\bottomrule
\end{tabular}
\end{table}

结论：偏移量在 $10^{-3}$--$10^{-1}\,$Hz 量级，远小于 EEG 峰频估计的典型精度
$\pm0.2$--$0.5\,$Hz。故 \cref{rem:offsetsmall} 的实用结论成立。

%-----------------------------------------------------------------------
\section{外周抖动的低通因子与量级估计}
\label{app:jitter}

设脉冲到达时刻为 $t_{n}=nT+\epsilon_{n}$，抖动
$\epsilon_{n}\sim\mathcal{N}(0,\sigma_{\tau}^{2})$ 独立同分布。第 $m=1$ 次
傅里叶分量的期望振幅按高斯特征函数衰减：
%
\begin{equation}
\bigl\langle e^{-\ii\Om\epsilon}\bigr\rangle
=\exp\!\left(-\frac{\sigma_{\tau}^{2}\Om^{2}}{2}\right)
\;\equiv\;\Xi_{\mathrm{jitter}} .
\end{equation}

\paragraph{在 $10\,$Hz 处的量级.}
$\Om=2\pi\times10=62.83\,\mathrm{rad\,s^{-1}}$，取上界
$\sigma_{\tau}=3\,$ms：
%
\begin{equation}
\frac{\sigma_{\tau}^{2}\Om^{2}}{2}
=\frac{(3\times10^{-3})^{2}\times(62.83)^{2}}{2}=0.0178 ,
\qquad
\Xi_{\mathrm{jitter}}=e^{-0.0178}=0.982 .
\end{equation}
%
即约 $1.8\%$ 衰减。

\paragraph{产生 $20\%$ 衰减所需的抖动.}
令 $\Xi_{\mathrm{jitter}}=0.8$，则 $\sigma_{\tau}^{2}\Om^{2}/2=0.223$，
%
\begin{equation}
\sigma_{\tau}=\frac{\sqrt{2\times0.223}}{62.83}=1.06\times10^{-2}\,\mathrm{s}
=10.6\,\mathrm{ms} .
\end{equation}
%
这远超外周神经传导抖动的生理范围。\textbf{结论：该机制在 mu 波段可忽略}，
仅在 gamma 波段（$\Om$ 大 $4$--$5$ 倍，指数中因子大 $16$--$25$ 倍）方案中
才需纳入。这是一个有信息量的负面结果：低频段外周通路是“干净”的传输
信道。\hfill$\blacksquare$

%-----------------------------------------------------------------------
\section{Rice 越阈公式到唯象门控}
\label{app:rice}

对零均值平稳高斯过程 $m(t)$，方差 $\sigma_{m}^{2}$，其对水平 $\Theta$ 的
向上穿越率由 Rice 公式给出：
%
\begin{equation}
\lambda(\Theta)=\frac{1}{2\pi}\frac{\sigma_{\dot m}}{\sigma_{m}}
\exp\!\left(-\frac{\Theta^{2}}{2\sigma_{m}^{2}}\right)
\equiv\nu_{0}\exp\!\left(-\frac{\Theta^{2}}{2\sigma_{m}^{2}}\right),
\end{equation}
%
即 \cref{eq:rice}。代入门控形式 \cref{eq:thetaeff}，$\Thf=\Theta_{0}+a\rho$：
%
\begin{align}
-\frac{\Thf^{2}}{2\sigma_{m}^{2}}
&=-\frac{\Theta_{0}^{2}+2a\Theta_{0}\rho+a^{2}\rho^{2}}{2\sigma_{m}^{2}}\notag\\
&=-\frac{\Theta_{0}^{2}}{2\sigma_{m}^{2}}
  -\frac{a\Theta_{0}}{\sigma_{m}^{2}}\rho
  -\frac{a^{2}}{2\sigma_{m}^{2}}\rho^{2} .
\end{align}
%
于是
%
\begin{equation}
\lambda=\underbrace{\nu_{0}e^{-\Theta_{0}^{2}/2\sigma_{m}^{2}}}_{\lambda_{0}}
\exp\!\left[-\underbrace{\frac{a\Theta_{0}}{\sigma_{m}^{2}}}_{\beta}\rho
            -\frac{a^{2}}{2\sigma_{m}^{2}}\rho^{2}\right].
\label{eq:app-lambda-full}
\end{equation}
%
截断二阶项即得 \cref{eq:lambda}。\hfill$\blacksquare$

\paragraph{二阶偏差的量级.}
二阶项相对一阶项之比为
%
\begin{equation}
\frac{a^{2}\rho^{2}/2\sigma_{m}^{2}}{a\Theta_{0}\rho/\sigma_{m}^{2}}
=\frac{a\rho}{2\Theta_{0}} .
\end{equation}
%
故当 mu 门控引起的阈值抬升 $a\rho$ 远小于基线阈值 $\Theta_{0}$ 时，一阶截断
成立。若实测剂量--反应曲线\emph{陡于}指数，即为二阶项的证据，可用于独立
估计 $a/\Theta_{0}$。

%-----------------------------------------------------------------------
\section{可辨识性：Fisher 信息矩阵的零方向}
\label{app:ident}

由 \cref{eq:Ei}，
%
\begin{equation}
E_i(\Om)=1-\exp\!\left[-\frac{B_i}{\sqrt{A_i^{2}+4(\wbar_i-\Om)^{2}}}\right],
\qquad
A_i=2\gamma_{\mathrm{eff},i}-K_i,\quad B_i=\beta_iF_i .
\end{equation}

\paragraph{$(\gamma_{\mathrm{eff}},K)$ 方向.}
由链式法则，
%
\begin{equation}
\frac{\partial E}{\partial\gamma_{\mathrm{eff}}}
=\frac{\partial E}{\partial A}\cdot
 \frac{\partial A}{\partial\gamma_{\mathrm{eff}}}
=2\,\frac{\partial E}{\partial A},
\qquad
\frac{\partial E}{\partial K}
=\frac{\partial E}{\partial A}\cdot\frac{\partial A}{\partial K}
=-\frac{\partial E}{\partial A} .
\end{equation}
%
故存在参数扰动方向
%
\begin{equation}
\bigl(\delta\gamma_{\mathrm{eff}},\,\delta K\bigr)\propto(1,\,2)
\quad\Longrightarrow\quad
\delta E=\frac{\partial E}{\partial\gamma_{\mathrm{eff}}}\cdot1
        +\frac{\partial E}{\partial K}\cdot2
        =2\frac{\partial E}{\partial A}-2\frac{\partial E}{\partial A}=0 ,
\end{equation}
%
即沿 $(1,2)$ 方向移动参数不改变任何观测量。Fisher 信息矩阵在该方向有严格
零特征值——$\gamma_{\mathrm{eff}}$ 与 $K$ \textbf{单独不可辨识}，仅组合
$A$ 可辨识。

\paragraph{$(\beta,F)$ 方向.}
$B=\beta F$ 为乘积，故 $\delta\ln B=\delta\ln\beta+\delta\ln F$。沿
%
\begin{equation}
\frac{\delta\beta}{\beta}=-\frac{\delta F}{F}
\end{equation}
%
移动同样不改变 $B$，故为第二个零方向。

\paragraph{结论.}
9 维参数空间在区制 I 下有效维数为 $3$，可辨识坐标即
\cref{eq:ident} 的 $\Theta_i=\{\wbar_i,A_i,B_i\}$。这不是模型缺陷，而是
\emph{正确识别了问题的内在自由度}：拟合应直接以 $(\wbar,A,B)$ 为参数，而非
在 9 维空间中做带正则化的病态优化。\hfill$\blacksquare$

\begin{remark}[对超临界区的补充]
区制 II 下 $\rho_{0}=\sqrt{1-2\gamma/K}$ 依赖 $\gamma/K$ 之\emph{比}而非差，
故该区制中可辨识组合与区制 I 不同。这意味着：\textbf{若能在同一被试上同时
测得亚临界型与超临界型响应特征，则 $\gamma$ 与 $K$ 可分别确定}。这为设计
参数分离实验提供了原则性依据。
\end{remark}

%-----------------------------------------------------------------------
\section{次谐波判别的候选实验设计}
\label{app:subharm}

本节给出 \cref{rem:onset} 所要求的、能把真次谐波锁定与起始瞬变／小波泄漏
区分开的实验设计要点。

\paragraph{混淆源清单.}
\begin{enumerate}[leftmargin=2.2em,itemsep=2pt]
\item \textbf{起始瞬变}：$1\,$s 脉冲串的 onset 为宽带事件，在 CWT 中沿低频
      展布，可产生虚假的 $\Om/2$ 功率。
\item \textbf{窗宽泄漏}：$1\,$s 矩形窗的频率分辨率约 $1\,$Hz；小波在低频端
      时间--频率权衡使 $5\,$Hz 附近估计方差大。
\item \textbf{包络周期}：$1\,$s on $/\,5\,$s off 的重复周期为 $1/6\,$Hz，其
      谐波族可与低频成分混淆。
\item \textbf{试次交替伪影}：若试次条件交替呈现，交替周期的一半可产生伪
      次谐波。
\end{enumerate}

\paragraph{判别设计.}
\begin{enumerate}[leftmargin=2.2em,itemsep=2pt]
\item \textbf{稳态段分析}：采用 $\ge5\,$s 连续刺激，剔除起始与结束各
      $500\,$ms，仅分析稳态段。真锁定应在稳态段持续存在；瞬变不应。
\item \textbf{相位一致性而非功率}：真 $\Om/2$ 锁定要求相邻两个刺激周期的
      响应相位\emph{交替}（周期倍化的定义特征）。计算 $\Om/2$ 处的
      试次间相位一致性并检验其与刺激相位的固定关系；瞬变无此结构。
\item \textbf{频率扫描}：若为真次谐波 tongue，其应随 $\Om$ 移动并保持
      $\Om/2$ 关系；若为伪影，则锚定在固定频率。
\item \textbf{振幅阈值}：高阶锁定的 tongue 通常比 $1{:}1$ 窄且需更高驱动。
      预测 $\Om/2$ 响应有\emph{强度阈值}，而线性瞬变随强度线性增长。
\end{enumerate}

\paragraph{判据.}
仅当 (1)(2)(3) 同时满足时，方可判定为非线性锁定证据。任一不满足，应归因于
\cref{rem:onset} 所述的平庸解释。


\bibliography{refs}

\end{document}
